\chapter{RSVP Fields and Action Principles}
\label{ch:rsvp-fields}

\section{Overview}

This chapter develops the dynamical and variational foundations of the
Relativistic Scalar--Vector Plenum (RSVP) field theory. Building on the semantic
manifold formalism of Chapter~\ref{ch:semantic-manifold}, we introduce the three
primary fields
\[
\Phi : \mathcal{M} \to \mathbb{R}, \qquad
\mathbf{v} \in \Gamma(T\mathcal{M}), \qquad
S : \mathcal{M} \to \mathbb{R}_{\ge 0},
\]
and construct the Lagrangian functional that governs their evolution. The
resulting Euler--Lagrange equations encode the dynamics of semantic
smoothing, directional flow, and entropy redistribution within the plenum.

RSVP differs from conventional physical field theories in one important
respect: the ``metric'' governing dynamical interactions is not assumed a
priori but emerges from the fields themselves. This reflects RSVP’s
interpretation of $\mathcal{M}$ as a \emph{semantic} rather than spacetime
manifold, where meaning gradients, uncertainty, and preferred flow directions
define the effective geometric structure.

\section{Field Content}

\subsection{The Scalar Potential Field \texorpdfstring{$\Phi$}{Φ}}

\begin{definition}[Semantic Potential]
A smooth scalar field $\Phi : \mathcal{M} \to \mathbb{R}$ represents the
\emph{semantic potential} of the plenum. The gradient $\nabla\Phi$ expresses
local tension or disequilibrium, and the RSVP dynamics tend to reduce these
gradients over time.
\end{definition}

Physically or semantically, $\Phi$ plays the role of a potential that drives
``smoothing forces.'' Steep gradients correspond to localized concentration of
structure, while shallow gradients represent near-equilibrium regions.

\subsection{The Flow Field \texorpdfstring{$\mathbf{v}$}{v}}

\begin{definition}[Anisotropic Flow]
A smooth vector field $\mathbf{v} \in \Gamma(T\mathcal{M})$ encodes preferred
directions along which semantic change propagates. Integral curves of
$\mathbf{v}$ represent anisotropic transport of semantic information.
\end{definition}

The field $\mathbf{v}$ biases the otherwise diffusive smoothing of $\Phi$,
introducing directed transport that reflects structural or contextual
constraints.

\subsection{The Entropy Density Field \texorpdfstring{$S$}{S}}

\begin{definition}[Entropy Density]
A smooth scalar field $S : \mathcal{M} \to \mathbb{R}_{\ge 0}$ represents local
entropy density, interpreted in RSVP as a measure of smoothness or uniformity
rather than disorder. Higher values of $S$ correspond to regions of greater
semantic regularity.
\end{definition}

Entropy in RSVP is not a measure of noise but of \emph{semantic coherence}.
Gradients of $S$ drive diffusion-like processes that push the plenum toward
uniformity.

\section{Effective Geometric Structure}

Because RSVP does not assume a background metric, we construct an
\emph{effective} pseudo-Riemannian structure $g$ from the fields themselves.
This ensures that the variational formulation has geometric meaning while
remaining consistent with the interpretation of $\mathcal{M}$ as a semantic
manifold.

\begin{definition}[Effective RSVP Metric]
Let $g$ be a symmetric $(0,2)$-tensor defined by
\[
g_{\mu\nu}
=
\alpha_1 \, \partial_\mu \Phi \, \partial_\nu \Phi
\;+\;
\alpha_2 \, S^{-1}\partial_\mu S \, S^{-1}\partial_\nu S
\;+\;
\alpha_3 \, v_\mu v_\nu ,
\]
where $\alpha_1,\alpha_2,\alpha_3$ are real constants that determine the
relative contribution of each field.
\]
\end{definition}

This metric reflects semantic geometry:
\begin{itemize}
\item gradients of $\Phi$ encode semantic tension;
\item gradients of $S$ encode smoothness stratification;
\item $\mathbf{v}$ introduces anisotropic directions.
\end{itemize}

We assume throughout that $g$ is nondegenerate on the strata where RSVP fields
vary.

The Levi--Civita connection of $g$ defines covariant derivatives, Laplacians,
and divergence operators used in the Lagrangian.

\section{Differential Operators}

\begin{definition}[Semantic Gradient and Divergence]
Given the effective metric $g$, define  
\[
\nabla_\mu \Phi = \partial_\mu \Phi, \qquad
\nabla^\mu \Phi = g^{\mu\nu} \nabla_\nu \Phi,
\]
and
\[
\nabla_\mu \mathbf{v}^\mu = \mathrm{div}(\mathbf{v}).
\]
\end{definition}

\begin{definition}[Semantic Laplacian]
For any scalar field $f$, define the Laplacian by
\[
\Delta f := \nabla_\mu \nabla^\mu f .
\]
\end{definition}

These operators appear in all subsequent RSVP equations.

\section{Construction of the RSVP Lagrangian}

The Lagrangian must satisfy several constraints:
\begin{enumerate}[label=(\roman*)]
\item invariance under diffeomorphisms of $\mathcal{M}$,
\item locality (dependence only on fields and finite jets),
\item compatibility with the semantic interpretation of smoothing and flow,
\item stability (boundedness of the action from below),
\item appropriate coupling to Polyxan curvature measures (see Ch.~\ref{ch:hypergraph-curvature}).
\end{enumerate}

We now define the RSVP Lagrangian density $\mathcal{L}$.

\subsection{Kinetic Terms}

The kinetic terms encode the smoothing tendency of the plenum:
\[
\mathcal{L}_{\mathrm{kin}}
=
\frac{1}{2} g^{\mu\nu} \partial_\mu \Phi \, \partial_\nu \Phi
\;+\;
\frac{\beta}{2} \, S^{-1} g^{\mu\nu} \partial_\mu S \, \partial_\nu S ,
\]
where $\beta > 0$ controls the relative smoothing rate of $S$.

\subsection{Flow Coupling}

The anisotropic term couples $\mathbf{v}$ to $\Phi$:
\[
\mathcal{L}_{\mathrm{flow}}
=
\gamma \, g^{\mu\nu} (\partial_\mu \Phi) v_\nu ,
\]
where $\gamma$ governs the strength of the directional bias.

\subsection{Potential Term}

The potential term encodes structural resistance, curvature penalties, and
contributions from Polyxan hypergraph geometry:
\[
\mathcal{L}_{\mathrm{pot}}
= - V(\Phi,S),
\]
where $V$ is assumed to be smooth, bounded below, and dependent only on the
local values of the fields.

\subsection{Complete Lagrangian}

\begin{definition}[RSVP Lagrangian]
The RSVP Lagrangian density is
\[
\mathcal{L}
=
\mathcal{L}_{\mathrm{kin}}
+
\mathcal{L}_{\mathrm{flow}}
+
\mathcal{L}_{\mathrm{pot}},
\]
and the action functional is
\[
\mathcal{S}[\Phi,\mathbf{v},S]
=
\int_{\mathcal{M}}
\sqrt{|g|} \, \mathcal{L} \, d^n x.
\]
\end{definition}

\section{Euler--Lagrange Field Equations}

We now compute the variations of $\mathcal{S}$ with respect to each field.

\subsection{Equation for \texorpdfstring{$\Phi$}{Φ}}

Variation with respect to $\Phi$ yields
\[
\nabla_\mu\nabla^\mu \Phi
\;+\;
\gamma \, \nabla_\mu v^\mu
\;-\;
\frac{\partial V}{\partial \Phi}
=
0.
\]

Interpretation:
\begin{itemize}
\item the Laplacian drives smoothing,
\item the divergence of $\mathbf{v}$ biases smoothing directionally,
\item the potential enforces structural constraints.
\end{itemize}

\subsection{Equation for \texorpdfstring{$S$}{S}}

\[
\beta \nabla_\mu (S^{-1} \nabla^\mu S)
\;-\;
\frac{\partial V}{\partial S}
=
0.
\]

Because of the $S^{-1}$ factor, this is a nonlinear diffusion equation; entropy
diffusion slows in regions of low smoothness.

\subsection{Equation for \texorpdfstring{$\mathbf{v}$}{v}}

Variation with respect to $v_\mu$ gives
\[
\gamma \nabla^\mu \Phi = 0
\quad \Rightarrow \quad
\nabla^\mu \Phi = 0 \ \text{if}\ \gamma\neq 0.
\]

Thus $\mathbf{v}$ aligns with the gradient of $\Phi$.

This alignment condition is the geometric core of RSVP: preferred directions of
semantic flow coincide with steepest semantic tension.

\section{Gradient Flow Formulation}

The RSVP dynamics can also be written as a gradient flow of the action:
\[
\partial_t \Phi = - \frac{\delta \mathcal{S}}{\delta \Phi},
\qquad
\partial_t S = - \frac{\delta \mathcal{S}}{\delta S},
\qquad
\partial_t \mathbf{v} = - \frac{\delta \mathcal{S}}{\delta \mathbf{v}}.
\]

This expresses RSVP evolution as an entropy-smoothing dynamical system on the
space of field configurations.

\section{Diffeomorphism Symmetry and BV Extension}

RSVP is diffeomorphism-invariant. The BV formalism is introduced here only
briefly; full treatment appears in Ch.~\ref{ch:rsvp-EL}.

\subsection{Gauge Transformations}

Infinitesimal diffeomorphisms generated by vector fields $\epsilon^\mu$ act as
\[
\delta_\epsilon \Phi = \epsilon^\mu \partial_\mu \Phi,
\qquad
\delta_\epsilon S = \epsilon^\mu \partial_\mu S,
\qquad
\delta_\epsilon v^\mu = \epsilon^\nu \partial_\nu v^\mu - v^\nu \partial_\nu \epsilon^\mu.
\]

\subsection{Ghosts and Antifields}

Introduce ghosts $c^\mu$ and antifields
\[
\Phi^*, \quad S^*, \quad v_\mu^*, \quad c_\mu^*.
\]

The extended action satisfies the classical master equation
\[
\{\mathcal{S}_{BV},\mathcal{S}_{BV}\} = 0.
\]

\section{Preview of Polyxan Coupling}

The coupling of RSVP fields to Polyxan hypergraphs involves a sheaf morphism
\[
\iota_x : \mathcal{A} \to J^1(\mathcal{M})_x,
\]
where $\mathcal{A}$ is the atom type space. RSVP fields regulate:
\begin{itemize}
\item node embeddings,
\item link curvature contributions,
\item hypergraph flow alignment along $\mathbf{v}$,
\item semantic tension redistribution via $\Phi$ and $S$.
\end{itemize}

Back-reaction occurs via curvature-dependent terms in the potential $V$.

This prepares the ground for Chapter~\ref{ch:rsvp-EL}, which derives global
stability results and the full Euler--Lagrange system.
